\section{Evaluation Setup}
\label{sec:experimental-setup}

\paragraph{Research Questions}\label{subsec:evaluation-questions}
Experiments first calibrate whether SRP accounts are usable
(calibration checks below); conditional on these checks, the analysis asks five
questions:
\begin{description}[leftmargin=!,labelwidth=2.7em]
    \item[RQ1:] What does SRP add to an ordinary logit-lens preference? That
    is, does splitting the scalar lens score into per-feature terms expose
    structure that the scalar alone hides?
    \item[RQ2:] How does the selected readout score change the sparse account?
    \item[RQ3:] Which readout features carry each side of token-vs-token
    preferences?
    \item[RQ4:] How does prompt context reweight readout-feature projections for
    a fixed token row? That is, how do the hidden-state projections $h^\top
    d_i$ shift across prompts while the token-row coefficients stay fixed?
    \item[RQ5:] Do the same readout-feature terms compose with two downstream methods? \emph{(a)} Does residual-stream attribution---which writes the final hidden state as a sum of per-component contributions from the embedding, attention heads, and MLP layers---yield per-component, per-feature accounting when paired with SRP? \emph{(b)} Do constrained interventions on selected decoder directions, fit on a discovery set of tokens, transfer to held-out tokens?
\end{description}

\paragraph{Models, Readout SAEs, and Held-Out Score Sets}\label{subsec:models-and-targets}
\emph{Model suite.}
The suite covers Qwen3.5 \citep{qwen35collection}, Gemma-4
\citep{gemma4official,gemma4collection}, Ministral \citep{liu2026ministral3},
and reasoning-distilled Qwen/Llama readouts (Qwen and Llama base checkpoints distilled on DeepSeek-R1 traces)
\citep{deepseekai2025deepseekr1}, spanning different tokenizers, LM-head row
distributions, parameter counts, and post-training recipes.
Model identifiers and details are in
Appendix~\ref{app:model-suite-selection}.

\emph{Qualitative display setting.}
Qualitative examples use one fixed low-error setting, Qwen3.5-2B
\(32\times\), \(k=256\), where \(D=32d_{\mathrm{model}}\) is dictionary width
and \(k\) is the TopK active-feature budget.

\emph{Dictionary selection and held-out score checks.}
We fit readout SAEs on final LM-head inputs, select them by held-out row
reconstruction, replacement fidelity, and sparse code usage, and reserve
held-out logit differences for downstream score-reconstruction checks.
Replacement uses 10,000 neutral text continuations from C4, the Colossal
Clean Crawled Corpus, per model
\citep{raffel2020exploring,dodge2021documenting}; score-reconstruction
evaluations use prebuilt prompt/score sets with fixed linear
coefficients, split assignment, and tokenization filters. These splits keep
dictionary selection separate from score-reconstruction evaluation.

\paragraph{Metrics, Baselines, and Operating Regimes}\label{subsec:metrics-and-baselines}
Calibration uses three families of checks, mirrored in the paragraph structure of \Cref{sec:results}: \emph{basis} (does the factorization preserve the LM head?), \emph{score} (does a selected logit, contrast, or margin survive the factorization?), and \emph{baselines} (does the learned sparse structure beat null and reference controls?).
\emph{Basis metrics.} Held-out row reconstruction and LM-head replacement on C4
hidden states check whether the readout SAE preserves the readout it replaces.
\emph{Score metrics.} Score reconstruction, sign preservation, and relative
error check whether a selected logit, contrast, or margin is reconstructed
accurately enough to read its per-feature contribution bars. \emph{Baselines.} Comparisons with nearest-row ridge (a dense low-rank reconstruction from nearest token rows in \(\WU\)), PCA (principal components at matched rank), shuffled sparse codes (permute the row--code assignment across tokens, breaking learned row--code structure while preserving sparsity), and random sparse patterns (Gaussian directions with random TopK supports, matching only the sparse capacity) check whether
learned token-row structure improves over dense low-rank structure, lexical
neighborhoods, or sparse capacity alone. Sign flips, large relative error,
weak replacement fidelity, or parity with baselines mark failure modes.

\emph{Scope and interpretation rule.}
All downstream analyses are readout-level: they describe how the LM head
decomposes a selected readout score, not how the model generates. We read feature bars
only when the selected readout score is reconstructed with small residual and, for
signed contrasts, the reconstructed and exact scores agree in sign. Feature
labels are descriptive summaries, while feature ids, signed contributions, and
residuals are the measured objects; the only claim about generation behavior
appears in the constrained-edit intervention in
\Cref{subsec:constrained-readout-edits}.
